\documentclass[border=2pt]{standalone}
\usepackage{amsmath, amsthm, amsfonts}
\usepackage{tikz-cd, kotex}
\usepackage{quiver}
\usepackage{Operators}
\usepackage{palette}
\usepackage{diagram-style}
\begin{document}

%%FIG 절댓값함수의 그래프 (Derivatives-1.svg)
\begin{tikzpicture}
\tikzset{
  axes/.style={->, black!50, line width=0.4pt},
  curve/.style={accent1, line width=0.9pt},
  dot/.style={fill=accent1, circle, inner sep=1.3pt}
}
\draw[axes] (-1.8,0) -- (1.8,0);
\draw[axes] (0,-0.35) -- (0,1.75);
\draw[curve] (-1.45,1.45) -- (0,0) -- (1.45,1.45);
\node[dot] at (0,0) {};
\end{tikzpicture}

%%FIG 세제곱근함수의 수직접선 (Derivatives-2.svg)
\begin{tikzpicture}
\tikzset{
  axes/.style={->, black!50, line width=0.4pt},
  curve/.style={accent1, line width=0.9pt},
  dot/.style={fill=accent1, circle, inner sep=1.3pt}
}
\draw[axes] (-2.0,0) -- (2.0,0);
\draw[axes] (0,-1.5) -- (0,1.5);
\draw[curve] plot[domain=-1.2:1.2, samples=160, smooth, variable=\t] ({\t*\t*\t},{\t});
\node[dot] at (0,0) {};
\end{tikzpicture}

%%FIG 2/3제곱함수의 첨점 (Derivatives-3.svg)
\begin{tikzpicture}
\tikzset{
  axes/.style={->, black!50, line width=0.4pt},
  curve/.style={accent1, line width=0.9pt},
  dot/.style={fill=accent1, circle, inner sep=1.3pt}
}
\draw[axes] (-2.1,0) -- (2.1,0);
\draw[axes] (0,-0.5) -- (0,1.9);
\draw[curve] plot[domain=-1.28:1.28, samples=200, variable=\t] ({\t*\t*\t},{\t*\t});
\node[dot] at (0,0) {};
\end{tikzpicture}

\end{document}
